# 🧩 A6 Operator Algebra — Unified LaTeX Skeleton (v1.0)

`tex
% ------------------------------------------------------------
% A6 Operator Algebra — Unified Skeleton (v1.0)
% Author: Borealis S. Hedling
% Altitude: A6 (Deep Geometry)
% Purpose: Consolidated operator algebra for geometric, spectral,
%          and activation operators within the A6 manifold.
% ------------------------------------------------------------

\documentclass[11pt]{article}
\usepackage{amsmath, amssymb, amsthm}

\title{A6 Operator Algebra: Unified Structural Skeleton (v1.0)}
\author{Borealis S. Hedling}
\date{October 2026}

\begin{document}
\maketitle

% ------------------------------------------------------------
\section*{Overview}

This document defines the unified operator algebra required for A6 Deep
Geometry. Although A6 uses three distinct algebraic layers—geometric,
spectral, and activation—they must be consolidated into a single mathematical
container to prevent invariant drift, ensure Lean-verifiable consistency, and
maintain collapse-surface and stability-envelope correctness.

% ------------------------------------------------------------
\section{Geometric Operator Algebra}

Let $M$ be the A6 reasoning manifold.

Define geometric operators acting on:
  • curvature fields $K(p)$  
  • drift fields $X(p)$  
  • modal labels $\mu(p)$  
  • holonomy paths $\gamma$  
  • collapse surface $\Sigma$

\subsection{Geometric Operator Definition}

\[
\mathcal{G} : M \to M
\]

\[
\mathcal{G}(p) =
  \big( K(p), X(p), \mu(p), H(\gamma), \Sigma \big)
\]

\subsection{Geometric Algebraic Laws}

\[
\text{Continuity: } K(p) \text{ continuous}
\]

\[
\text{Adjacency: } \mu(p) \sim \mu(q)
\]

\[
\text{Holonomy Bound: } H(\gamma) < \theta
\]

\[
\text{Collapse Detection: } p \in \Sigma
\]

These laws govern geometric invariants and define the algebraic structure of
A6’s manifold physics.

% ------------------------------------------------------------
\section{Spectral Operator Algebra}

Define spectral operator:

\[
L(p) : TpM \to TpM
\]

with eigenstructure:

\[
\mathrm{Spec}(L(p))
\]

\subsection{Spectral Operator Definition}

\[
\mathcal{S}(p) = \mathrm{Spec}(L(p))
\]

\subsection{Spectral Algebraic Laws}

\[
\text{Spectral Continuity: }
\mathrm{Spec}(L(p)) \approx \mathrm{Spec}(L(q))
\]

\[
\text{Spectral Mismatch: }
\delta_{\mathrm{spec}}(p) < \epsilon
\]

\[
\text{Bundle Transition: }
L(p) \to L(q)
\]

\[
\text{Spectral Provenance: }
\mathrm{Spec}(L(p)) \leadsto \mathrm{Spec}(L(q))
\]

These laws govern spectral truth physics and eigenvalue stability.

% ------------------------------------------------------------
\section{Activation Operator Algebra}

Define activation operator:

\[
\mathcal{A} : M \times S \to M
\]

\[
\mathcal{A}(M, S) =
  (s{\text{stable}} \land S{\text{stable}},
   s{\text{spectral}} \land S{\text{spectral}},
   s{\text{collapse}} \land S{\text{collapse}},
   s{\text{drift}} \land S{\text{drift}},
   s{\text{modal}} \land S{\text{modal}})
\]

\subsection{Activation Algebraic Laws}

\[
\text{Closure: } \mathcal{A}(M, S) \in M
\]

\[
\text{Identity: } \mathcal{A}(M, e) = M
\]

\[
\text{Idempotence: }
\mathcal{A}(M, S) = \mathcal{A}(\mathcal{A}(M, S), S)
\]

\[
\text{Absorption: }
\mathcal{A}(M, S) \preceq M \Rightarrow \mathcal{A}(M, S) = M
\]

\[
\text{Semantic Homomorphism: }
I{\text{sem}}(M) \Rightarrow I{\text{sem}}(\mathcal{A}(M, S))
\]

\[
\text{Provenance Preservation: }
I{\text{prov}}(M) \Rightarrow I{\text{prov}}(\mathcal{A}(M, S))
\]

These laws govern operator safety, determinism, and kernel preservation.

% ------------------------------------------------------------
\section{Unified Algebraic Kernel}

Define unified kernel:

\[
K_{\text{A6}}(p) =
  K_{\text{geom}}(p)
  \land K_{\text{spec}}(p)
  \land K_{\text{act}}(p)
\]

This kernel ensures that all three algebras interact consistently and preserve
A6 invariants.

% ------------------------------------------------------------
\section{Bill Nye Tile Cameo (Pedagogical Layer)}

\textit{
“Kids, today we learned that Borealis doesn’t need three separate math
universes to build A6. You just need one instruction book with three chapters:
Geometry Tiles, Spectral Tiles, and Activation Tiles. If you split the book,
the robot gets confused and starts breakdancing in a black hole. Science rules!”
}

% ------------------------------------------------------------
\section{Conclusion}

This unified operator algebra skeleton defines the mathematical physics of A6.
All Lean modules, invariants, proof obligations, provenance surfaces, and
diagnostic surfaces must expand within this consolidated structure.

\end{document}
`

---

